\documentclass[10pt]{article}

\usepackage[utf8]{inputenc}

\usepackage[T1]{fontenc}

\usepackage{amsmath}

\usepackage{amsfonts}

\usepackage{amssymb}

\usepackage[version=4]{mhchem}

\usepackage{stmaryrd}

\usepackage{mathrsfs}

\usepackage{bbold}



\begin{document}

$M^{d}=\left(M_{t}^{d}, \mathscr{F}_{t}\right)$ - чисто разрывный локальный мартингал, к-рый может быть записан в виде



\[

M_{t}^{d}=\int_{0}^{t} \int_{|x| \leqslant 1} x d(\mu-v),

\]



где $d \mu=\mu(\omega, d t, d x)$ - случайная мера скачков семимартингала $X$, то есть



$\mu(\omega,(0, t], \Gamma)=\sum_{0<s \leqslant t} I\left(\Delta X_{s} \in \Gamma\right), \Gamma \in \mathscr{B}(\mathbb{R} \backslash\{0\})$,



a $d v=v(\omega, d t, d x)$ - ее компенсатор. Поскольку



\[

\sum_{0<s \leqslant t} \Delta X_{s} I\left(\left|\Delta X_{s}\right|>1\right)=\int_{0}^{t} \int_{|x|>1} x d \mu

\]



то всякий семимартингал $X$ допускает представление $X_{t}=X_{0}+B_{t}+M_{t}^{c}+\int_{0}^{t} \int_{|x| \leqslant 1} x d(\mu-v)+\int_{0}^{t} \int_{|x|>1} x d \mu$, к-рое наз. к а но ни ч е с ким пред с т а в лен и е м (р а з ло же ни е м).



Набор (предсказуемых) характеристик $\left(B,\left\langle M^{c}\right\rangle, v\right)$, где $\left\langle M^{c}\right\rangle$ - квадратич. характеристика $M^{c}$, т. е. такой предсказуемый возрастающий процесс, что $\left(M^{c}\right)^{2}-\left\langle M^{c}\right\rangle$ является локальным мартингалом, наз. т р и п л е т о м локальных (предсказуемых) характеристик $X$.



Jum.: [i] J a c o d J., Calcul stochastique et problemès de mar-tingales, B., 1979 (Lecture notes in mathematics, № 714).



СЕПАРАБЕЛЬНАЯ АЛГЕБРА - конечномерная полупростая ассоциативная алгебра $A$ над полем $k$, остающаяся полупростой при любом расширении $K$ поля $k$ (т. е. алгебра $A \otimes_{k} K$ полупроста для любого поля $K \supseteq k)$. Алгебра $A$ сепарабельна тогда и только тогда, когда центры простых компонент этой алгебры (см. Ассоциативные кольца $u$ алгебры) являются сепарабельными расширениями поля $k$.



Лит.: [1] В а н Д е р В а р д е н Б. Л., Алгебра, пер. с нем., 2 изд., М., 1979; [2] К ә р т и с Ч., Р а Й н е р И., Теория представлений конечных групп и ассоциативных алгебр, пер. с англ., М., 1969. СЕIАРӒБЕЛЬНОЕ ОТОБРАЖЕНИЕ - домИнантный морфизм $f$ неприводимых алгебраич. многообразий $X$ и $Y, f: X \rightarrow Y$, для к-рого поле $K(X)$ является сепарабельным расширением подполя $f^{*} K(Y)$ (изоморфного $K(Y)$ ввиду доминантности). Несепарабельные отображения существуют только тогда, когда характеристика $p$ основного поля больше нуля. Если $f$ - конечный морфизм и его степень не делится на $p$, то он сепарабелен. При С. о. для точек в нек-ром открытом подмножестве $U \subset X$ диффференциал $(d f)_{x}$ отображения $f$ сюръективно отображает касательное пространство $T_{X, x}$ в $T_{Y, f(x)}$ и наоборот: если точки $x$ и $f(x)$ неособые и $(d f)_{x}$ - сюръективен, то $f$ есть С. о.



Термин «сепарабельность» употребляется для морфизмов и в ином смысле. Морфизм $f: X \rightarrow Y$ схем $X$ и $Y$ наз. с е п а р а 6 е л ь ным, или (чаще) о т д ел и м ы м, если диагональ в $X \times_{Y} X$ замкнута. Композиция отделимых морфизмов отделима; $f: X \rightarrow Y$ отделим тогда и только тогда, когда для любой точки $y \in Y$ существует такая окрестность $V \ni y$, что морфизм $f: f^{-1}(V) \rightarrow V$ отделим. Морфизм аффинных схем всегда отделим. Существует критерий отделимости для нётеровых схем.



СЕПАРАБЕЛЬНОЕ ПРОСТРАНСТВО - топологическое пространство, обладающее счетной базой. Про такие пространства иногда говорят, что они удовлетворяют второй аксиоме счетности. М. И. Войцеховский.



СЕПАРАБЕЛЬНОЕ РАСІІИРЕНИЕ п $о$ Л Я - pacширение $K / k$ такое, что для нек-рого натурального $n$ поля $K$ и $k^{p^{-n}}$ линейно разделены над $k$ (см. Линейно разделенные расширения). Расширение, не являющееся сепарабельным, наз. н е с е п а р а б е л ь вы м. В дальнейпем рассматриваются только алгебраич. расширения (о трансцендентных сепарабельных расширениях см. Трансцендентное расширение). Конечное расширение сепарабельно тогда и только тогда, когда отображение следа $\mathrm{Sp}: K \rightarrow k$ является ненулевой функцией. Алгебраич. расширение сепарабельно, если любое конечное его подрасширение сепарабельно. В характеристике 0 все расширения сепарабельны.



C. р. образуют о т м е ч е ны й к л а с с р а cши р е н и й, т. е. в башне полей $L \supset K \supset k$ расширение $L / k$ сепарабельно тогда и только тогда, когда сепарабельны и $L / K$ и $K / k$. Если $K_{1} / k$ и $K_{2} / k$ суть С. р., то и $K_{1} K_{2} / k$ сепарабельно; для С. р. $K / k$ и произвольного расширения $L / k$ расширение $K L / L$ снова сепарабельно. Расширение $K / k$ сепарабельно тогда и только тогда, когда оно допускает погружение в некpoe расширение Галуа $L / k$. При этом для конечного расширения $K / k$ число различных $k$-изоморфизмов поля $K$ в $L$ совпадает со степенью $[K: k]$. Любое конечное С. р. является простым.



Многочлен $f(x) \in k[X]$ наз. с е п а р а б е л ь ны м над $k$, если его неприводимые множители не имеют кратных корней. Алгебраич. әлемент $\alpha$ наз. с е п ар а б е л ь ны м (над $k$ ), если он является корнем сепарабельного над $k$ многочлена. В противном случае $\alpha$ наз. н е с е пі а р а б е л ь н ы м. Элемент $\alpha$ наз. ч и с т о н е с е п а р а б е л ь ны м над $k$, если $\alpha^{p^{n}} \in k$ для нек-рого $n$. Неприводимый многочлен $f(x)$ несепарабелен тогда и только тогда, когда производная $f^{\prime}(x)$ тождественно равна 0 (это возможно только в случае, когда $k$ имеет характеристику $p$ и $f(x)=$ $=f_{1}\left(x^{p}\right)$. Произвольный многочлен $f(x)$ однозначно представим в виде $f(x)=g\left(x p^{e}\right)$, где $g(x)$ - сепарабельный многочлен. Степень многочлена $g(x)$ и число $e$ наз. соответственно р е д у ц и р о в а н н о й $\quad$ т епе н ь ю и и н д е к с о м многочлена $f(x)$.



Пусть $L / k$ - произвольное алгебраич. расширение. Все элементы поля $L$, сепарабельные над $k$, образуют поле $K$, к-рое является максимальным С. р. поля $k$, содержащимся в $L$. Поле $K$ наз. с е п а р а б е л ьны м за мы к а н и е м поля $k$ в $L$. Степень $[K: k]$ наз. с е п а р а б ө льн о й с те пень ю расширения $L / k$, а степень $[L: K]-$ н е с е п а р а б е л ьной с т е п е вь ю, или с тепенью не с е п ар а б е ль н о с т и. Несепарабельная степень равна нек-рой степени числа $p=\operatorname{char} k$. Если $K=k$, то поле $k$ наз. с е па р а бельн о з а мк нуты т в $L$. В этом случае распииение $L / k$ наз. ч и с т о н е с еп а р а б е л в ны м. Расширение $K / k$ чисто несепарабельно тогда и только тогда, когда



\[

K \subset k^{p^{-\infty}}=\bigcup_{n} k^{p^{-n}}

\]



т. е. когда любой элемент поля $K$ чисто несепарабелен над $k$. Чисто несепарабельные распирения поля $k$ образуют отмеченный класс расширений. Если расширение $K / k$ одновременно сепарабельно и чисто несепарабельно, то $K=k$.



Лит. см. при ст. Расширение поля. Л. В. Кузъмин.



СЕПАРАБЕЛЬНЫЙ ПРОЦЕСС - случайный процесс, поведение траекторий к-рого по существу определяется их поведением на нек-ром счетном пространстве. Именно, определенный на полном вероятностном пространстве $\{\Omega, \mathscr{F}, \mathrm{P}\}$ действительный случайный процесс $\left\{X_{t}, t \in T\right\}$, где $T$ - подмножество действительной прямой $\mathbb{R}$, с е п а р а б е ле н о т н о с ит е л ь н о к л а с с а $\mathcal{A}$ подмножеств $\mathbb{R}$, если существует счетное множество $T_{1} \subset T$ (с е п а р а н т а) и множество $N \subset \mathcal{F}, \mathrm{P}(N)=0$, такое, что для любого $A \subset \mathcal{A}$ и любого открытого интервала $I \subset \mathbb{R}$



\[

\bigcap_{t \in I T_{1}}\left\{X_{t} \in A\right\} \backslash \cap_{t \in I T}\left\{X_{t} \in A\right\} \subset N

\]





\end{document}